\section*{Chapitre \ref{ch:hf:news-and-event-trading} --- Trading d'actualités et d'événements}

\begin{solution}{exo:hf:news-and-event-trading:1}
$10\times3+10\times2+10\times1=60$ lots de tick.
\end{solution}

\begin{solution}{exo:hf:news-and-event-trading:2}
$(1-p)\times2=p\times10$~: $p=1/6\approx16.7~\%$.
\end{solution}

\begin{solution}{exo:hf:news-and-event-trading:3}
719~\$~; $238\times12.50\times12=\$35\,700$ par an pour les publications d'un seul contrat.
\end{solution}

\begin{solution}{exo:hf:news-and-event-trading:4}
Les cotations périmées sont limitées et prises dans l'ordre de latence~: la capacité du premier palier en couvre l'essentiel avant l'arrivée du second, et les fournisseurs
mettent à jour leurs cotations avant même l'arrivée des paliers lents.
\end{solution}

\begin{solution}{exo:hf:news-and-event-trading:5}
Dans les premières secondes, le mark-out dépasse la fourchette (trois ticks contre un)~; après trente secondes, le flux supplémentaire s'est en grande partie dissipé, si bien qu'un teneur de marché tardif
en manque la meilleure part.
\end{solution}

\begin{solution}{exo:hf:news-and-event-trading:6}
Lorsque quelques traders informés négocient pendant deux secondes, les prix se sont déplacés vers la nouvelle valeur avant la publication publique~; les autres traders faisaient alors face à un prix qui
reflétait déjà l'information, et la publication publique levait moins d'incertitude. Un trading précoce concentré, selon eux, réduisait
l'avantage de la vitesse après la publication publique.
\end{solution}

\begin{solution}{exo:hf:news-and-event-trading:7}
Seul le palier de 0.05 milliseconde arrive avant la mise à jour des fournisseurs~: il prend tout~; la valeur cédée passe de 223 à 197 lots de tick par publication
(la capacité du premier palier devient contraignante) et la perte de notre teneur de marché passe de 57.5 à 52.7.
\end{solution}

\begin{solution}{exo:hf:news-and-event-trading:8}
Les horodatages des fournisseurs indiquent le moment où le fournisseur a envoyé le chiffre, et non celui où le trader l'a reçu~; le backtest suppose une latence que la firme n'a pas, et
exécute les ordres aux cotations que des traders plus rapides ont prises en premier. Utiliser les horodatages de réception sur le site du trader et la profondeur au repos après les paliers plus rapides.
\end{solution}

\begin{solution}{pb:hf:news-and-event-trading:1}
\begin{enumerate}
\item Voir la \cref{def:hf:news-and-event-trading:lockup} et la \cref{def:hf:news-and-event-trading:race}.
\item L'accès anticipé de deux secondes (2007--juin 2013) a produit un trading concentré et une formation des prix en moins de 200 millisecondes~; la publication par paliers a pu
réduire l'avantage des traders les plus rapides.
\item Thomson Reuters a mis fin à l'accès anticipé de deux secondes (juillet 2013)~; un tweet piraté de l'Associated Press a fait bouger les indices américains pendant quelques minutes (avril 2013)~; Musk et
Tesla ont versé chacun 20~millions de dollars au titre des tweets sur un financement sécurisé (septembre 2018).
\item Une fois assez rapide, une mauvaise lecture coûte plus que le fait d'arriver en second.
\item $J=\mathrm{round}(3z)$ ticks~; dix niveaux de cinquante lots périmés plus les nôtres~; quatre paliers de 200 lots à 0.05, 1, 5, 50 ms~; les fournisseurs se mettent à jour après 20
ms.
\item 88.1~\%, 11.6~\%, 0.3~\% et 0.
\item Ses cotations rejoignent les niveaux après celles des autres~; les preneurs consomment les files d'attente par l'avant.
\item 57.5 lots de tick, 719~\$.
\item Voir la \cref{def:hf:news-and-event-trading:protocol}.
\item 872, 911, 1\,052 et 599 lots de tick.
\item En restant~: $872-58=814$~; en retirant puis en réinscrivant ses cotations à cinq secondes~: 1\,052.
\item Des limites de \omterm{def:hf:depositary-receipts-and-dual-listings:prerelease}{pré-émission}, des contrôles de vraisemblance sur les chiffres analysés, une taille maximale par événement, un kill switch accessible, et l'annulation automatique elle-même.
\item 88~\% au palier le plus rapide, 12~\% au suivant~; le teneur de marché perd 58 lots de tick en restant et rien avec son protocole, et gagne 1\,052 contre 814
sur deux minutes.
\item Les fournisseurs qui laissent des cotations périmées et, via des fourchettes plus larges autour des publications, tout le monde.
\item Moins de valeur cédée, et la totalité allant aux tout premiers.
\item Avec des déclencheurs (volatilité, signaux liés aux titres) qui retirent automatiquement les cotations, et des règles de réentrée.
\item La course se joue au moment de réception chez le trader.
\item Réaction aux titres des réseaux sociaux.
\item Mark-outs des exécutions par seconde après chaque publication, moyennés sur de nombreuses publications.
\item Éviter les pertes de la course et être présent lorsque le flux est important mais n'est plus toxique.
\end{enumerate}
\end{solution}

\begin{solution}{iq:hf:news-and-event-trading:1}
Retirer ou élargir automatiquement ses cotations quelques secondes avant~; réentrer à une fourchette plus large une fois passés les mark-outs des premières secondes, puis
le resserrer à mesure qu'ils s'atténuent.
\iqlookfor{retrait automatisé et réentrée mesurée.}
\end{solution}

\begin{solution}{iq:hf:news-and-event-trading:2}
Feed handler, analyseur, contrôle de vraisemblance, mise en correspondance avec les instruments, génération d'ordres, contrôles de risque, passerelle. Défaillances~: retard du flux, erreur d'analyse, un
chiffre hors de sa plage, prévision consensuelle périmée, contrôle de risque qui bloque ou ne bloque pas.
\iqlookfor{la chaîne et ses points de défaillance.}
\end{solution}

\begin{solution}{iq:hf:news-and-event-trading:3}
Régresser le mouvement de chaque instrument pendant les premières secondes sur la surprise statistique standardisée des publications passées~; tester hors échantillon et par régime.
\iqlookfor{estimation et validation hors échantillon.}
\end{solution}

\begin{solution}{iq:hf:news-and-event-trading:4}
Dimensionner la taille pour que la perte de queue (dix fois le gain moyen) sur un événement reste dans la limite de perte quotidienne~; poser un cap sur le nombre d'événements simultanés.
\iqlookfor{dimensionnement sur la queue de distribution.}
\end{solution}

\begin{solution}{iq:hf:news-and-event-trading:5}
L'authenticité de la source, la confirmation par l'entreprise ou par une agence de presse, si les négociations ont été suspendues, et le mouvement de prix observé jusqu'ici.
\iqlookfor{vérification avant action.}
\end{solution}

\begin{solution}{iq:hf:news-and-event-trading:6}
Les gains à partir de $t$ sont $\int_t^H\lambda(1+Be^{-s/T_1})(h-a e^{-s/T_2})\,ds$~; la dérivée en $t$ est l'opposé de l'intégrande en $t$, donc l'optimum est
où $h=a e^{-t/T_2}$~: $t^\ast=T_2\ln(a/h)$, soit ici $5\ln3=5.5$ secondes.
\iqlookfor{réentrer lorsque l'exécution marginale atteint le seuil de rentabilité.}
\end{solution}
